\chapter*{Introduction g\'en\'erale}
\addcontentsline{toc}{chapter}{Introduction g\'en\'erale}
\markboth{Introduction g\'en\'erale}{Introduction g\'en\'erale}

Les objets connect\'es produisent des communications nombreuses et
h\'et\'erog\`enes dont l'analyse manuelle devient rapidement impraticable. La
d\'etection d'intrusions fond\'ee sur les flux r\'eseau offre une observation
moins intrusive que l'inspection du contenu, mais elle impose trois exigences:
un sch\'ema de donn\'ees stable, une \'evaluation scientifique reproductible et
une cha\^ine de traitement capable de conserver les preuves utiles à
l'investigation.

Ce projet d'\'et\'e \'etudie la question suivante: comment construire un
prototype reproductible qui classe des observations de flux IoT, signale les
comportements suspects et pr\'esente à l'op\'erateur les informations n\'ecessaires
à leur examen, sans confondre un r\'esultat obtenu sur RT-IoT2022 avec une
validation en environnement de production? Le travail couvre la pr\'eparation
contr\^ol\'ee du jeu de donn\'ees, l'\'evaluation d'une cascade supervis\'ee,
l'ingestion durable, la persistance des r\'esultats et une interface
d'investigation. La collecte NFStream en direct et l'efficacit\'e sur un r\'eseau
cible restent hors du p\'erim\`etre valid\'e.

Le rapport est organis\'e en quatre chapitres. Le premier \'etablit le cadre
scientifique, la probl\'ematique et la d\'emarche de recherche. Le deuxi\`eme
formalise les besoins et les cas d'utilisation. Le troisi\`eme pr\'esente la
conception dynamique, statique et architecturale. Le quatri\`eme d\'ecrit
l'environnement de r\'ealisation, les principaux choix d'impl\'ementation et les
interfaces essentielles du prototype.
